% Dirichlet Round 11 English TeX
% Work: G. Lejeune Dirichlet, Werke II, pp. 159--181.
% English translation checked against Werke II scan images 172--194.
\documentclass[11pt]{article}
\usepackage[a4paper,margin=1in]{geometry}
\usepackage[T1]{fontenc}
\usepackage[utf8]{inputenc}
\usepackage{lmodern}
\usepackage{amsmath,amssymb}
\usepackage{microtype}
\setlength{\parindent}{1.2em}
\setlength{\parskip}{0.25em}
\begin{document}

\begin{center}
{\Large G. Lejeune Dirichlet's Works. II.}\par\vspace{1.2em}
{\Large XIV.}\par\vspace{0.5em}
{\LARGE \textbf{Simplification of the Theory\ of Binary Quadratic Forms\ of Positive Determinant.}}\par\vspace{1em}
{\large G. Lejeune Dirichlet}\par\vspace{0.5em}
{\small Liouville, Journal de Mathématiques pures et appliquées, Series II, Volume II, pp. 353--376.}
\end{center}

\begin{center}
{\large \textbf{Simplification of the Theory of Binary Quadratic Forms\ of Positive Determinant.}}\par
{\small Read at the Academy of Sciences on 13 July 1854.\footnote{The notice in the Academy report for 1854, p. 384, reads: ``Mr. Dirichlet read: On a property of continued fractions and its use for simplifying the theory of quadratic forms of positive determinant.'' K.}}
\end{center}

The more the higher arithmetic has gained in scope through Gauss's epoch-making work and through later investigations, the more desirable it becomes that access to this beautiful branch of analysis should be made as easy as possible by simplifying its elementary part. With this aim I have, in several earlier papers, prefaced my investigations with new proofs of the known propositions required for them. The present paper has a similar simplification in view, and is devoted to the theory of quadratic forms of positive determinant. In its previous shape this theory requires, as is well known, quite detailed discussions; the following treatment shows how they may be removed.

I begin with some remarks on continued fractions. Although their essential content is not new, they must be put in the form needed for the application made here.

\section*{\S. 1.}

A finite or infinite continued fraction such as
\[
  \alpha+\frac{1}{\beta+\frac{1}{\gamma+\cdots}}
\]
will be denoted below by
\[
  (\alpha,\beta,\gamma,\ldots).
\]
We shall consider only continued fractions all of whose terms are integers, except of course the last term in the case where the expansion has not been carried to an end and where that term, as a complete quotient, may have any value. Of special importance for arithmetical investigations are the continued fractions whose terms are all positive except possibly the first, for which the value zero is also allowed. A positive irrational number $\omega$ can be expressed in this form in only one way. We call this expression of $\omega$, or, when $\omega$ is negative, the expression of its absolute value with a preceding minus sign, the normal continued-fraction expansion of $\omega$.

We first have to derive the normal expansion of an irrational number $\omega$ from a continued fraction of the form
\[
 \omega=(\alpha,\beta,\ldots,\mu,\nu,p,q,r,\ldots,u,v,\ldots),
\]
in which all terms are positive only from $p$ onward, including $p$ itself. This can be done by a sequence of transformations which leave unchanged all terms following a sufficiently remote $u$. The number of new terms which finally replace $\alpha,\beta,\ldots,\nu$ differs from the number of these old terms by an even or an odd number according as $\omega$ is positive or negative.

First suppose that $\nu$ is not the first term. Then $\mu,\nu$ and some immediately following terms may be replaced, all other terms being left unchanged, by new terms whose number differs by an even number, and all of which are positive except possibly the first, which may be zero or negative. The cases are as follows. If $\nu=0$, the three terms $\mu,0,p$ are replaced by the single term $\mu+p$. If $\nu=-n$, one distinguishes
\[
 n>1,
 \qquad n=1,
\ p>1,
 \qquad n=1,
\ p=1,
\]
and uses, respectively,
\[
 (\mu,-n,p,q,\ldots)=(\mu-1,1,n-2,1,p-1,q,\ldots),
\]
\[
 (\mu,-1,p,q,\ldots)=(\mu-2,1,p-2,q,\ldots),
\]
\[
 (\mu,-1,1,q,r,s,\ldots)=(\mu-q-2,1,r-1,s,\ldots).
\]
The change in the number of terms is respectively $2,0,-2$. If one of the non-negative differences
\[
 n-2,
\quad p-1,
\quad p-2,
\quad r-1
\]
becomes zero, the zero and the two adjacent positive terms are replaced by a single term equal to their sum.\footnote{Lagrange already observed that negative terms can be removed from a continued fraction; but the equation he gave for this purpose, which is the first of the equations above, is insufficient, because when $p=1$ it introduces a new negative term. Applying the same equation again returns one to the original continued fraction. Oeuvres de Lagrange, t. 2 p. 624.}

By repeating this procedure one makes all terms from the second onward positive. If the first term is then not negative, the operation is finished. If instead the first term is $-a$, so that
\[
 \omega=(-a,b,c,d,\ldots),
\]
then, according as $b>1$ or $b=1$, one writes
\[
 \omega=-(a-1,1,b-1,c,\ldots),
 \qquad\hbox{or}\qquad
 \omega=-(a-1,c+1,d,\ldots).
\]
This again gives the result just asserted.

\section*{\S. 2.}

I. Let two quantities $\omega,\Omega$ and integers $\alpha,\beta,\gamma,\delta$, with $\alpha\ne0$, satisfy
\[
 \omega=\frac{\gamma+\delta\Omega}{\alpha+\beta\Omega},
 \qquad
 \alpha\delta-\beta\gamma=1.
\]
Then one can always form an equation
\[
 \omega=(\lambda,m,\ldots,r,\sigma,
\Omega)
\]
in which, among the integers $\lambda,m,\ldots,r,\sigma$, only the first and the last can be zero or negative; the intermediate terms, if any occur, are positive and are even in number.

Since the signs of $\alpha,\beta,\gamma,\delta$ may all be changed at the same time, suppose $\alpha>0$. If $\alpha=1$, then immediately
\[
 \omega=\frac{\gamma+(\beta\gamma+1)\Omega}{1+\beta\Omega}
       =(\gamma,\beta,\Omega).
\]
If $\alpha>1$, expand $\gamma/\alpha$ in the ordinary way, choosing all remainders positive:
\[
 \frac{\gamma}{\alpha}=(\lambda,m,\ldots,r).
\]
Only $\lambda$ can be zero or negative, and the number of terms $m,
\ldots,r$ may be assumed even, since an end term $r>1$ may, if necessary, be replaced by $(r-1,1)$. Let the convergents be
\[
 \frac{\lambda}{1},
 \quad
 \frac{\lambda m+1}{m},
 \quad\ldots,
 \quad
 \frac{\varphi}{f},
 \quad
 \frac{\gamma}{\alpha}.
\]
The last has denominator $\alpha$, and the standard relation is
\[
 \alpha\varphi-
\gamma f=1.
\]
Comparison with $\alpha\delta-
\beta\gamma=1$ gives
\[
 \beta=\alpha\sigma+f,
 \qquad
 \delta=\gamma\sigma+\varphi
\]
with an integer $\sigma$. Thus $\delta/\beta$ is annexed to the sequence of convergents by the new term $\sigma$, and
\[
 \omega=(\lambda,m,\ldots,r,\sigma,
\Omega).
\]

II. We shall need the special case in which $\alpha,\beta,\gamma,\delta$ are all positive and also
\[
 \gamma\geq\alpha,
 \qquad
 \delta>\gamma.
\]
Then $\lambda$ and $\sigma$ are positive. If $\alpha=1$, this is immediate, for then $\lambda=\gamma$ and $\sigma=\beta$. If $\alpha>1$, then $\lambda$ is the integer immediately below $\gamma/\alpha$, hence positive. Also $\sigma>0$; for from
\[
 \delta=\gamma\sigma+\varphi
\]
one would obtain $\delta=\varphi<\gamma$ if $\sigma=0$, and a negative $\delta$ if $\sigma<0$. Denoting the positive numbers $\lambda,
\sigma$ by $l,s$, we have in this special case
\[
 \frac{\gamma}{\alpha}=(l,m,\ldots,r),
 \qquad
 \frac{\delta}{\beta}=(l,m,\ldots,r,s),
 \qquad
 \omega=(l,m,\ldots,r,s,
\Omega),
\]
where all terms $l,m,
\ldots,r,s$ are positive and even in number.

\section*{\S. 3.}

We now turn to the quadratic forms
\[
 ax^2+2bxy+cy^2=(a,b,c),
\]
all of which are to have the same positive determinant
\[
 D=b^2-ac.
\]
The positive integer $D$ is arbitrary except that it must not be a square. Hence the outer coefficients $a,c$ are never zero, and once $D$, the middle coefficient, and one of the two outer coefficients are given, the other outer coefficient, and therefore the form itself, is determined.

To each form $(a,b,c)$ we attach the equation
\[
 a+2b\omega+c\omega^2=0,
\]
whose roots are written
\[
 \omega=\frac{-b\mp\sqrt D}{c}
\]
with the unchanged third coefficient $c$ as denominator. According to the upper or lower sign, these are called the first and second roots of the form. A form is completely determined by its determinant and either one of its roots: if the same value is a first root, or a second root, of both $(a,b,c)$ and $(A,B,C)$, then
\[
 \frac{-b\mp\sqrt D}{c}=\frac{-B\mp\sqrt D}{C},
\]
and the irrationality of $\sqrt D$ gives $B=b$, $C=c$.

Two forms
\[
 ax^2+2bxy+cy^2,
 \qquad
 AX^2+2BXY+CY^2
\]
will be called equivalent only in the proper sense: there exists a substitution
\[
 x=\alpha X+\beta Y,
 \qquad
 y=\gamma X+\delta Y,
\]
with
\[
 \alpha\delta-
\beta\gamma=1
\]
which transforms the first form into the second. Substitutions of determinant $-1$ are excluded throughout.

I. If two forms are equivalent, then their roots of the same name, $\omega$ and $\Omega$, and the coefficients of the substitution satisfy
\[
 \omega=\frac{\gamma+\delta\Omega}{\alpha+\beta\Omega}.
\]
Equivalently,
\[
 \Omega=\frac{\gamma-
\alpha\omega}{\beta\omega-
\delta}.
\]
After substituting the value of $\omega$ and rationalising the denominator, the right hand side becomes
\[
 \frac{-M\mp\sqrt D}{N},
\]
where
\[
 M=a\alpha\beta+b(\alpha\delta+\beta\gamma)+c\gamma\delta,
 \qquad
 N=a\beta^2+2b\beta\delta+c\delta^2.
\]
These are precisely the second and third coefficients $B,C$ of the transformed form.

II. Conversely, if this equation holds for a pair of roots of the same name and if $\alpha,\beta,\gamma,\delta$ satisfy $\alpha\delta-\beta\gamma=1$, then the two forms are equivalent and the indicated substitution carries the first into the second. For one obtains
\[
 \frac{-B\mp\sqrt D}{C}=\frac{-M\mp\sqrt D}{N},
\]
hence $B=M$, $C=N$.

III. We shall often use neighbouring forms, that is, forms
\[
 (a,b,a'),
 \qquad
 (a',b',a'')
\]
such that
\[
 b+b'\equiv0\pmod {a'}.
\]
They are always equivalent. Applying to the first form the substitution
\[
 \begin{pmatrix}0&1\\-1&\delta\end{pmatrix}
\]
one obtains first coefficient $a'$ and second coefficient $-b-a'\delta$, which equals $b'$ when
\[
 \delta=-\frac{b+b'}{a'}.
\]
For the corresponding roots
\[
 \omega=\delta-
\frac1{\omega'},
 \qquad
 \omega'=-\frac1{\omega-
\delta}.
\]

\section*{\S. 4.}

A form $(a,b,c)$ is called reduced if, of the two roots
\[
 \frac{-b-
\sqrt D}{c},
 \qquad
 \frac{-b+
\sqrt D}{c},
\]
the first has absolute value greater than $1$, the second has absolute value less than $1$, and the two roots have opposite signs. Then $b>0$ and $b<\sqrt D$; hence $-ac=D-b^2>0$, so $a$ and $c$ have opposite signs. The sign of the first root is the sign of $a$ and is opposite to that of $c$.

If $(a,b,c)$ is reduced, so is $(c,b,a)$; each root of one is the reciprocal of the opposite-name root of the other.

There are only finitely many reduced forms of a given determinant $D$. They are obtained by taking each positive $b<\sqrt D$, finding all positive and negative factors $c$ of $D-b^2$ whose absolute values lie between $\sqrt D+b$ and $\sqrt D-b$, and then putting
\[
 a=-\frac{D-b^2}{c}.
\]

Let $(a,b,a')$ be a reduced form. Among its right neighbours $(a',b',a'')$ there is exactly one reduced form. If the neighbour is to be reduced, its first root $\omega'$ must have sign opposite to the first root $\omega$ of $(a,b,a')$. In
\[
 \omega'=-\frac1{\omega-
\delta},
\]
choose $\delta$ so that $\omega'$ is an improper fraction and has sign opposite to $\omega$. Equivalently, $\omega-
\delta$ must be a proper fraction with the same sign as $\omega$. There is exactly one such integer: the integer adjacent to $\omega$ on the same side of $\omega$ as zero. It is nonzero, has the sign of $\omega$, and its absolute value is immediately below $|\omega|$.

This neighbour is indeed reduced. If $\omega$ now denotes the second root in the same equation, then $\omega'$ has the sign of $-\delta$; also $\omega-
\delta$ is improper, hence $\omega'$ is proper and has the sign required for the second root.

For the middle coefficient of this unique right neighbour set, for the first root,
\[
 \omega-
\delta=\sigma,
\]
where $\sigma$ is a proper fraction with the same sign as $\omega$. From
\[
 b'=-b-a'\delta
\]
we get
\[
 b'=\sqrt D+a'\sigma.
\]
Since $\sigma$ has sign opposite to $a'$, $b'$ lies between $\sqrt D$ and $\sqrt D\mp a'$, with the upper or lower sign according as $a'$ is positive or negative. Together with
\[
 b'\equiv -b\pmod {a'},
\]
this determines $b'$ without ambiguity. Similarly there is a unique reduced left neighbour $({}'a,{}'b,a)$.

\section*{\S. 5.}

Starting from a reduced form $\varphi_0$, let $\varphi_1$ be the reduced form immediately to its right, let $\varphi_2$ be the right neighbour of $\varphi_1$, and so on; do the same to the left. One obtains an infinite two-sided sequence of equivalent reduced forms
\[
 \ldots,\varphi_{-2},\varphi_{-1},\varphi_0,
 \varphi_1,\varphi_2,\ldots.
\]
Because only finitely many reduced forms have the same determinant, forms in this sequence repeat. Since the first coefficients alternate in sign, identity can occur only at an even difference of indices. If two forms are identical, every pair of forms standing the same distance to the same side of them is also identical.

Let $\varphi_{2n}$ be the first form after $\varphi_0$ identical with $\varphi_0$. Then
\[
 \varphi_0,\varphi_1,\ldots,
\varphi_{2n-1}
\]
are all distinct and form a period. Thus
\[
 \varphi_\mu=\varphi_\nu
 \quad\Longleftrightarrow\quad
 \mu\equiv\nu\pmod {2n}.
\]
The same criterion holds for the equality of corresponding roots $\omega_\mu,\omega_\nu$.

Let $\delta_\nu$ be the greatest integer contained in the absolute value of the first root $\omega_\nu$, taken with the sign of $\omega_\nu$. Then
\[
 \omega_\nu=\delta_\nu-
\frac1{\omega_{\nu+1}}.
\]
The congruence of indices implies equality of the $\delta$'s, though not conversely.

Choose the index so that forms of even index have positive first coefficient. Then the signs of $\omega_\nu$ and $\delta_\nu$ are $(-1)^\nu$. Write
\[
 \delta_\nu=(-1)^\nu k_\nu,
\]
with $k_\nu>0$. The numbers $k_\nu$ have period $2n$, and
\[
 \omega_\nu=(-1)^\nu(k_\nu,k_{\nu+1},\ldots,k_{2n+
u-1},k_\nu,k_{\nu+1},\ldots).
\]
For the second root $\omega'_\nu$,
\[
 \frac1{\omega'_\nu}=(-1)^{\nu-1}(k_{\nu-1},k_{\nu-2},\ldots,k_{\nu-2n},k_{\nu-1},k_{\nu-2},\ldots).
\]
Thus the period for the reciprocal of the second root is the reverse of the period for the first root.

The period $k_0,k_1,\ldots,k_{2n-1}$ has no shorter even period. Otherwise $\omega_0$ and $\omega_{2m}$, hence also $\varphi_0$ and $\varphi_{2m}$, would coincide with $2m<2n$.

All reduced forms of a given determinant can be partitioned into periods of this sort: form a period from any reduced form, and, if any reduced forms remain, repeat the operation until none are left.

\section*{\S. 6.}

We now decide whether two given forms are equivalent. Since every form can be reduced and forms in the same period are equivalent, it remains only to decide whether forms in different periods can be equivalent.

Choose in each period a form with positive first coefficient:
\[
 \varphi_0=(a,b,c),
 \qquad
 \Phi_0=(A,B,C),
\]
and write their first roots in normal continued fractions
\[
 \omega_0=(k_0,k_1,k_2,\ldots),
 \qquad
 \Omega_0=(K_0,K_1,K_2,\ldots).
\]
Suppose the forms are equivalent, and that a substitution of determinant $1$ takes the first into the second. Then
\[
 \omega_0=\frac{\gamma+\delta\Omega_0}{\alpha+\beta\Omega_0}.
\]
Here $\alpha$ cannot be zero: otherwise $\gamma=\pm1$ and $A=c$, contrary to the sign assumptions. Hence by \S. 2
\[
 \omega_0=(\lambda,m,\ldots,r,\sigma,
 K_0,K_1,K_2,\ldots).
\]
Let the number of inserted terms be $2g$. When this continued fraction is brought to normal form, the number of terms up to a sufficiently remote unchanged term $K_\nu$ changes, since $\omega_0>0$, by an even number $2h$. Thus for all sufficiently large $\nu$,
\[
 K_\nu=k_{2g+2h+
u}.
\]
Replacing $\nu$ by $2Ni+\nu$, where $2Ni$ is a large multiple of the period $2N$ of the $K$-sequence, and reducing $2g+2h+2Ni$ modulo $2n$ to $2m$, gives for all $\nu\ge0$
\[
 K_\nu=k_{2m+
u}.
\]
Therefore
\[
 \Omega_0=\omega_{2m},
 \qquad
 \Phi_0=\varphi_{2m}.
\]
The second form lies in the period of the first. Thus forms from different periods cannot be equivalent.

\section*{\S. 7.}

It remains to indicate how the rest of the theory is modified in the points that depend on the proof just given.

The operations which decide equivalence give one substitution taking one form to the other. It remains to derive all such substitutions from one of them. This reduces to describing all substitutions which carry a form into itself. We may suppose that the coefficients of the form are coprime, since removing a common factor does not change these substitutions.

Let $(a,b,c)$ be primitive, and let $\sigma$ be the positive greatest common divisor of $a,2b,c$. Then $\sigma$ is $1$ or $2$, and the latter occurs only when $D=b^2-ac$ has the form $4h+1$. All substitutions carrying the form into itself are given by
\[
 \alpha=\frac{t-bu}{\sigma},
 \qquad
 \beta=-\frac{cu}{\sigma},
 \qquad
 \gamma=\frac{au}{\sigma},
 \qquad
 \delta=\frac{t+bu}{\sigma},
\]
where $t,u$ run through all integer solutions of
\[
 t^2-Du^2=\sigma^2.
\]
The complete solution of this indeterminate equation is easily obtained from its solution in the least positive numbers.

For a reduced form the transformation problem can be solved directly. Suppose $a>0$ and consider substitutions whose four coefficients are positive. Let
\[
 \omega=(k_0,k_1,\ldots,k_{2n-1},k_0,k_1,\ldots)
\]
be the normal periodic continued fraction for the first root. If $\gamma/\alpha$ and $\delta/\beta$ are successive convergents, the second corresponding to the last term $k_{2n-1}$ of a period, then
\[
 \alpha\delta-\beta\gamma=1,
 \qquad
 \omega=\frac{\gamma+\delta\omega}{\alpha+\beta\omega},
\]
and hence the corresponding substitution carries the form into itself.

Conversely, let $\alpha,\beta,\gamma,\delta$ be the coefficients of such a positive substitution. Then
\[
 \beta\omega^2+(\alpha-
\delta)\omega-
\gamma=0.
\]
Since one root lies between $1$ and $\infty$ and the other between $-1$ and $0$, the left hand side is negative at $\omega=1$ and positive at $\omega=-1$. Consequently
\[
 \gamma-
\alpha>\beta-
\delta,
 \qquad
 \delta-
\gamma>\alpha-
\beta,
\]
and hence
\[
 \gamma\geq\alpha,
 \qquad
 \delta>\gamma.
\]
All hypotheses of \S. 2, II, are therefore satisfied:
\[
 \frac{\gamma}{\alpha}=(l,m,\ldots,r),
 \qquad
 \frac{\delta}{\beta}=(l,m,\ldots,r,s),
 \qquad
 \omega=(l,m,\ldots,r,s,
\omega).
\]
Substitution of the periodic expansion of $\omega$ shows that $l,m,\ldots,r,s$ consists of one or more complete periods, and that the two fractions are successive convergents ending at the end of a period. The smallest positive substitution corresponds to the end of the first period.

The link with the Pell-type equation is as follows. From the formulas for $\alpha$ and $\delta$,
\[
 \alpha\delta=\frac{t^2-b^2u^2}{\sigma^2}
 =\frac{t^2-Du^2}{\sigma^2}-\frac{ac u^2}{\sigma^2}
 =1-
\frac{ac u^2}{\sigma^2}.
\]
Since $-ac>0$, $\alpha$ and $\delta$ have the same sign, which is the sign of $t$ because
\[
 \alpha+
\delta=\frac{2t}{\sigma}.
\]
Similarly, $\beta$ and $\gamma$, unless both zero, have the sign of $u$. Thus the positive substitutions just considered arise from positive $t,u$; because $\beta$ and $\gamma$ increase with $u$, the substitution expressed in the smallest numbers gives the least positive values of $t,u$, namely
\[
 t=\frac{\sigma}{2}(\delta+
\alpha),
 \qquad
 u=\frac{\sigma}{2b}(\delta-
\alpha).
\]


\begin{center}
\textbf{Addition to this paper, by the author.}
\end{center}

In order that the reader need not turn back to the writings cited in the note concerning \S. 7, we shall prove in a few words the formulas recalled at the point to which that note refers. The necessary and sufficient conditions for the substitution
\[
  \begin{pmatrix}\alpha&\beta\\ \gamma&\delta\end{pmatrix}
\]
to transform the form $(a,b,c)$ into itself are plainly expressed by
\[
  \alpha\delta-\beta\gamma=1,
  \qquad
  a(\alpha^2-1)+2b\alpha\gamma+c\gamma^2=0,
  \qquad
  a\alpha\beta+2b\beta\gamma+c\gamma\delta=0,
\]
the third of these having been simplified by replacing $\alpha\delta$ with $\beta\gamma+1$. Adding the last two equations after multiplying them respectively by $\delta$ and $-\gamma$, and then after multiplying them by $-\beta$ and $\alpha$, one obtains the equivalent system
\[
  \alpha\delta-\beta\gamma=1,
  \qquad
  a(\alpha-\delta)+2b\gamma=0,
  \qquad
  a\beta+c\gamma=0.
\]
Since $a\ne0$, if we put $\gamma=a\psi$, the rational quantity $\psi$ is completely determined by $\gamma$, and the last two equalities may be replaced by
\[
  \delta-\alpha=2b\psi,
  \qquad
  \beta=-c\psi,
  \qquad
  \gamma=a\psi.
\]
Now $\delta-\alpha$, $\beta$, and $\gamma$ are integers, and if the greatest common divisor of $a,2b,c$ is denoted by $\sigma$, $\psi$ must necessarily be of the form $u/\sigma$, where $u$ is an integer. Thus
\[
  \delta-\alpha=\frac{2bu}{\sigma},\qquad
  \beta=-\frac{cu}{\sigma},\qquad
  \gamma=\frac{au}{\sigma}.
\]
Substitution in
\[
  (\delta+\alpha)^2=(\delta-\alpha)^2+4\alpha\delta
       =(\delta-\alpha)^2+4\beta\gamma+4
\]
and multiplication by $\sigma^2$ give
\[
  [\sigma(\delta+\alpha)]^2=4[(b^2-ac)u^2+\sigma^2]
      =4(Du^2+\sigma^2).
\]
It follows that $\sigma(\delta+\alpha)$ is even. Writing
\[
  \sigma(\delta+\alpha)=2t,
\]
with $t$ an integer, the last equation takes the form
\[
  t^2-Du^2=\sigma^2,
\]
and the integers $\alpha,\beta,\gamma,\delta$ are expressed by
\[
 \alpha=\frac{t-bu}{\sigma},\qquad
 \beta=-\frac{cu}{\sigma},\qquad
 \gamma=\frac{au}{\sigma},\qquad
 \delta=\frac{t+bu}{\sigma}.
\]
Conversely, every integral solution $(t,u)$ of this equation gives integral values of $\alpha,\beta,\gamma,\delta$, and these values satisfy the equations that express the sufficient conditions for the corresponding substitution to transform $(a,b,c)$ into itself. The integrality is immediate when $\sigma=1$; when $\sigma=2$ it follows from the fact that $D$ and $b$ are odd in that case, so that $t$ and $u$ are both even or both odd. The remaining assertion is verified by direct substitution. \hfill(Toul, 14 August 1857.)

\begin{center}
\textbf{Extract from a letter of Dirichlet to Liouville.}
\end{center}

On looking over the paper recently translated by Hoüel, I noticed that it could be made more complete without lengthening it, by using a remark that I have had occasion to make in my course. The last two paragraphs of \S. 4 are to be replaced by the following.

To determine the coefficient $b'$, observe that the reciprocal value of the first of the roots belonging to the form $(a',b',a'')$ must be a quantity $\sigma$ whose absolute value is less than unity and whose sign is that of $a'$. Hence
\[
  \sigma=\frac{-b'+\sqrt D}{a'},\qquad
  b'=\sqrt D-a'\sigma.
\]
Thus $b'$ lies between $\sqrt D$ and $\sqrt D-(a')$, where $(a')$ denotes the absolute value of $a'$. This condition, together with
\[
  b'\equiv-b\pmod {a'},
\]
determines $b'$ uniquely and easily.

One proves in exactly the same way that there always exists a unique reduced form contiguous to the proposed one on the left; this also follows from the preceding case by the earlier observation.

We finish the paragraph by showing that one can always obtain a reduced form equivalent to any given form $(a,b,a')$. To this end form the sequence
\[
 (a,b,a'),\quad (a',b',a''),\quad (a'',b'',a'''),\ldots,
\]
each form being obtained from the preceding one by the process just indicated. The forms so obtained will eventually all be reduced, and this will occur at the latest when one reaches a form whose first coefficient does not exceed the third in absolute value. That point cannot fail to occur, since the positive integers
\[
 (a),\ (a'),\ (a''),\ldots
\]
cannot decrease indefinitely.

It remains only to prove that a form
\[
  (A,B,C),
\]
with $(A)\leq(C)$ and with $(B)$ between $\sqrt D$ and $\sqrt D-(A)$, is reduced. By the last condition,
\[
  B=\sqrt D-A\sigma,
  \qquad
  \sigma=\frac{-B+\sqrt D}{A},
  \qquad
  \frac1\sigma=\frac{-B-\sqrt D}{C},
\]
where $\sigma$ has absolute value less than one and the sign of $A$. From the last two equations and from $(A)\leq(C)$ it follows that the roots
\[
  \frac{-B-\sqrt D}{C},\qquad
  \frac{-B+\sqrt D}{C}
\]
of the form have, in absolute value, the first greater than one and the second less than one. Since the first root is thus numerically greater than the second, $B$ is necessarily positive; then the next-to-last equation, in which $\sigma$ and $A$ have the same sign, gives $B<\sqrt D$. Consequently the two roots have opposite signs. Q.E.D.


\end{document}
